% ==================================================================
% Conscious Point Physics:
% The Eight Gluons as hDP Structures on the Tetrahedral 600-Cell Subgraph
% Strong Sector Series #3 — Version 1
% Companion to: SS#1 (overview), SS#2 (SU(3) algebra);
%               EW Series #1–5 v3; C14 (confinement), C15 (color charge)
% GitHub: CPP/series_strong/cpp_ss3_gluons_v1.tex
% ==================================================================

\documentclass[11pt,a4paper]{article}
\usepackage[margin=1in]{geometry}
\usepackage[utf8]{inputenc}
\usepackage[T1]{fontenc}
\usepackage{amsmath,amssymb,amsthm}
\usepackage{graphicx}
\usepackage{svg}
\svgsetup{inkscapelatex=false,inkscapeformat=pdf,inkscapepath=svgdir}
\usepackage{caption,float,booktabs,parskip,xcolor}
\usepackage[numbers,sort&compress]{natbib}
\usepackage{hyperref}
\hypersetup{colorlinks=true,linkcolor=blue,urlcolor=cyan,citecolor=red}

\newtheorem{theorem}{Theorem}[section]
\newtheorem{proposition}[theorem]{Proposition}
\newtheorem{corollary}[theorem]{Corollary}
\newtheorem{remark}[theorem]{Remark}
\newtheorem{openprob}[theorem]{Open Problem}

\newcommand{\lP}{l_P}
\newcommand{\SSV}{\mathrm{SSV}}
\newcommand{\hDP}{\mathrm{hDP}}
\newcommand{\qCP}{\mathrm{qCP}}
\newcommand{\phig}{\varphi}
\newcommand{\ket}[1]{|#1\rangle}
\newcommand{\bra}[1]{\langle#1|}
\newcommand{\as}{\alpha_s}

\title{\textbf{Conscious Point Physics:\\
The Eight Gluons as hDP Structures\\
on the Tetrahedral 600-Cell Subgraph}\\[6pt]
{\large Strong Sector Series \#3 --- Version 1}}

\author{Thomas Lee Abshier, ND and Grok (x.AI)\\
Hyperphysics Institute\\
\url{https://hyperphysics.com}\\
\texttt{drthomas007@protonmail.com}}

\date{23 March 2026}

\begin{document}
\maketitle

\begin{abstract}
In Conscious Point Physics (CPP), the eight QCD gluons emerge as
transient hybrid dipole pair (hDP) configurations on the three edges
of the qCP tetrahedral base $\{V_1, V_2, V_3\}$.  Six
\emph{color-changing} gluons correspond to hDP pairs propagating
along the three base edges, and two \emph{color-neutral diagonal}
gluons correspond to phase-difference modes between base vertices.
This paper derives four gluon properties directly from the
tetrahedral cage geometry:

\textbf{Masslessness:} Gluons are transient open-path hDP pairs,
not stable closed polyhedral structures.  The absence of a
confining closed subgraph means zero SS~Vector compression energy
and hence zero mass --- exact parallel to the photon ($\lambda=0$
DP-Sea mode).

\textbf{Spin-1:} Each tetrahedral base edge has a definite
orientation vector in 3D space.  The hDP emitted along edge
$V_i \to V_j$ inherits this direction as its polarisation axis,
giving vector (spin-1) character.

\textbf{Color octet:} The 8 gluons transform in the adjoint
representation of SU(3), with Casimir $C_A = 3$ (verified
numerically).  The color charge carried by a gluon is the
difference between the vertex labels of its emission and
absorption endpoints.

\textbf{Non-Abelian self-coupling:} The 3-gluon vertex emerges
from nested tetrahedral hoppings: when two gluons act
sequentially, the antisymmetric part of the combined operator
$T^aT^b - T^bT^a = if^{abc}T^c$ gives the cubic term in the
Yang-Mills Lagrangian.  The 4-gluon vertex follows similarly.

The Casimir invariants $C_F = 4/3$ (fundamental) and $T_F = 1/2$
(Dynkin index) are derived from the operator algebra proved in
SS\#2 \citep{cpp_ss2}.  Color confinement forces color-neutral
hadrons: baryons (one quark at each base vertex) and mesons
(vertex + complementary antivertex).
\end{abstract}

\tableofcontents
\newpage

%=====================================================================
\section{Introduction}
%=====================================================================

SS\#2 \citep{cpp_ss2} proved that the eight tetrahedral hopping
operators $T^a$ equal the standard SU(3) generators $T^a =
\lambda^a/2$ exactly.  This paper asks: what are these operators
\emph{physically}?

Each $T^a$ generates a transition between two color states, or a
phase difference between them.  In QCD, the mediators of these
transitions are gluons.  In CPP, they are hDP pairs emitted along
specific edges of the tetrahedral base, carrying the color-label
change from emission vertex to absorption vertex.

The key structural distinction from the electroweak bosons (W, Z, H)
is topology:
\begin{itemize}
\item EW bosons are \emph{closed} structures --- bracelet, icosahedron,
  dodecahedron --- with SS~Vector compression energy and hence mass.
\item Gluons are \emph{open} transient paths --- hDP pairs in transit
  along a single edge --- with no closed confinement geometry and
  hence zero mass.
\item The photon is the EW analogue: the $\lambda = 0$ DP-Sea mode,
  massless for the same structural reason (no stable closed subgraph).
\end{itemize}

%=====================================================================
\section{The Eight hDP Configurations}
%=====================================================================

\subsection{Color-changing gluons (six)}

The tetrahedral base has three undirected edges, each connecting
a pair of base vertices:
$\{V_1, V_2\}$, $\{V_1, V_3\}$, $\{V_2, V_3\}$.
Each directed edge $V_i \to V_j$ supports a color-changing hDP
transition $\ket{i} \to \ket{j}$.  The real and imaginary
combinations give two gluon states per undirected edge:

\begin{align}
g_{r\bar{g}} &= \ket{r}\!\bra{g}
= T^1 + iT^2,\quad&
g_{g\bar{r}} &= \ket{g}\!\bra{r}
= T^1 - iT^2, \label{eq:g12}\\
g_{r\bar{b}} &= \ket{r}\!\bra{b}
= T^4 + iT^5,\quad&
g_{b\bar{r}} &= \ket{b}\!\bra{r}
= T^4 - iT^5, \label{eq:g13}\\
g_{g\bar{b}} &= \ket{g}\!\bra{b}
= T^6 + iT^7,\quad&
g_{b\bar{g}} &= \ket{b}\!\bra{g}
= T^6 - iT^7. \label{eq:g23}
\end{align}

Each gluon $g_{i\bar{j}}$ carries color $i$ and anticolor $\bar{j}$,
changing the quark's color from $j$ to $i$ upon absorption.  This
is the CPP description of what QCD calls a ``color-anticolor'' gluon
in the adjoint (octet) representation.

\subsection{Color-neutral diagonal gluons (two)}

Two diagonal operators capture phase differences between base
vertices without changing color:
\begin{align}
g_3 &= T^3 = \frac{1}{2}(\ket{r}\!\bra{r} - \ket{g}\!\bra{g}),
\label{eq:g3} \\
g_8 &= T^8 = \frac{1}{2\sqrt{3}}
(\ket{r}\!\bra{r} + \ket{g}\!\bra{g} - 2\ket{b}\!\bra{b}).
\label{eq:g8}
\end{align}

These are the two neutral gluons in QCD.  In CPP, $g_3$ measures
the red--green phase asymmetry and $g_8$ the red+green vs.\ blue
asymmetry.  Both are ``cage breathing modes'' --- oscillations
in the relative phase of the three base vertices without any
vertex-to-vertex hopping.

\subsection{Complete gluon table}

\begin{table}[H]
\centering
\caption{The 8 gluons: CPP hDP structure, operator, and color content.}
\label{tab:gluons}
\begin{tabular}{@{}cllll@{}}
\toprule
Gluon & Notation & CPP hDP process & Operator & Color content \\
\midrule
1 & $g_{r\bar{g}}$ & $V_2 \to V_1$ real & $T^1 + iT^2$ & Carries $r\bar{g}$ \\
2 & $g_{g\bar{r}}$ & $V_1 \to V_2$ real & $T^1 - iT^2$ & Carries $g\bar{r}$ \\
3 & $g_{r\bar{b}}$ & $V_3 \to V_1$ real & $T^4 + iT^5$ & Carries $r\bar{b}$ \\
4 & $g_{b\bar{r}}$ & $V_1 \to V_3$ real & $T^4 - iT^5$ & Carries $b\bar{r}$ \\
5 & $g_{g\bar{b}}$ & $V_3 \to V_2$ real & $T^6 + iT^7$ & Carries $g\bar{b}$ \\
6 & $g_{b\bar{g}}$ & $V_2 \to V_3$ real & $T^6 - iT^7$ & Carries $b\bar{g}$ \\
7 & $g_3$ & $V_1$--$V_2$ phase diff & $T^3$ & Neutral ($r\bar{r} - g\bar{g}$) \\
8 & $g_8$ & $(V_1\!+\!V_2)$--$V_3$ phase & $T^8$ & Neutral ($r\bar{r}+g\bar{g}-2b\bar{b}$) \\
\bottomrule
\end{tabular}
\end{table}

%=====================================================================
\section{Masslessness from Open-Path Topology}
%=====================================================================

\begin{theorem}[Gluon masslessness]
\label{thm:massless}
Every gluon has zero SS~Vector compression energy and hence zero
rest mass.
\end{theorem}

\begin{proof}
The SS~Vector compression energy of a CPP structure is
(EW Series \#2 eq.~4 \citep{cpp_ew2}):
\begin{equation}
E_{\rm conf}
= f_{\rm geom} \times \text{sea\_strength}
  \times \frac{\hbar c}{\lP^3}
  \times 4\pi r_{\rm eff}
  \times \phig^{-3} \times \eta.
\label{eq:econf}
\end{equation}
This energy is nonzero only for \emph{closed} polyhedral
structures (bracelet, icosahedron, dodecahedron, fullerene) because
only closed structures produce a net inward SS~Vector gradient that
confines DI bits.  An open-path hDP pair propagating along a single
tetrahedral edge has no bounding closed subgraph: the bits flow
freely along the edge without any confinement geometry.  Therefore
$f_{\rm geom} = 0$ for every gluon configuration, giving
$E_{\rm conf} = 0 = m_g c^2$.\qed
\end{proof}

\begin{remark}[Gluon vs.\ photon masslessness: same mechanism]
The photon is the $\lambda = 0$ mode of the 600-cell vertex
adjacency spectrum: no confinement energy, neutral, massless
(EW Series \#1 \citep{cpp_ew1}).  Gluons are the analogous
open-path modes of the \emph{tetrahedral cell} structure:
transient hDP pairs with no closed geometry, carrying color
instead of electric charge but massless for identical reasons.
The unification is complete: \emph{all massless gauge bosons in
the SM are open-path DP-Sea modes in CPP.}
\end{remark}

\begin{remark}[Contrast with massive EW bosons]
The W, Z, and Higgs acquire mass precisely because they are closed
polyhedral structures:
\begin{center}
\begin{tabular}{@{}llcl@{}}
\toprule
Boson & Topology & $f_{\rm geom}$ & Mass \\
\midrule
W & Closed bracelet & $0.219$ & $80.4$~GeV \\
Z & Closed icosahedron & $0.263$ & $91.2$~GeV \\
H & Closed dodecahedron & $0.1415$ & $125.1$~GeV \\
$g$ & Open edge (transient) & $0$ & $0$ \\
$\gamma$ & $\lambda=0$ DP-Sea mode & $0$ & $0$ \\
\bottomrule
\end{tabular}
\end{center}
\end{remark}

%=====================================================================
\section{Spin-1 from Edge Orientation}
%=====================================================================

\begin{theorem}[Gluon spin]
\label{thm:spin1}
Each gluon carries spin angular momentum $J = 1$ (vector boson).
\end{theorem}

\begin{proof}[Argument]
After stereographic projection from the 600-cell's 4D coordinates
to 3D (via $\mathbf{r}_{3D} = (x,y,z)/(1-w)$), each tetrahedral
base edge $V_i \to V_j$ maps to a definite oriented line segment
$\hat{e}_{ij}$ in 3D space.

A DI-bit hDP emitted along $\hat{e}_{ij}$ carries:
\begin{itemize}
\item A momentum $\mathbf{p} \parallel \hat{e}_{ij}$ (propagation
  along the edge direction).
\item An angular momentum $\mathbf{L} = \mathbf{r} \times \mathbf{p}$
  from the directed edge, giving one unit of angular momentum
  $|\mathbf{L}| = \hbar$ in the direction perpendicular to
  $\hat{e}_{ij}$.
\end{itemize}
One unit of angular momentum $\hbar$ is spin-1.  The two
polarisation states of a massless spin-1 boson correspond to the
two transverse directions relative to $\hat{e}_{ij}$, matching the
two physical polarisations of the gluon.\qed
\end{proof}

\begin{remark}[Same mechanism as the photon]
The photon acquires spin-1 from the direction of DI-bit propagation
in the DP~Sea (SR/GR companion C1 \citep{abshier_sr}).  The gluon
acquires spin-1 from the direction of the tetrahedral cage edge.
In both cases, spin is geometrically set by the propagation
direction, not by an intrinsic angular momentum property of the
quanta.
\end{remark}

%=====================================================================
\section{Color Octet Representation and Casimir Invariants}
%=====================================================================

\subsection{Adjoint (octet) representation}

\begin{theorem}[Gluons transform in the adjoint representation]
\label{thm:adjoint}
Under SU(3)$_c$ color rotations, the 8 gluon fields $A^a_\mu$
transform in the adjoint (8-dimensional) representation:
\begin{equation}
A^a_\mu \to A^a_\mu + f^{abc}\theta^b A^c_\mu,
\label{eq:adj_transform}
\end{equation}
where $f^{abc}$ are the structure constants derived in SS\#2
\citep{cpp_ss2}.
\end{theorem}

\begin{proof}
By definition, the adjoint representation generators are
$(T^a_{\rm adj})_{bc} = -if^{abc}$.  Since the structure
constants $f^{abc}$ are fully antisymmetric and satisfy the
Jacobi identity (SS\#2 Theorem~3 \citep{cpp_ss2}), the adjoint
matrices form a valid 8-dimensional representation of su(3).
The transformation law~\eqref{eq:adj_transform} is then the
standard infinitesimal adjoint action.\qed
\end{proof}

\subsection{Casimir invariants}

\begin{theorem}[Casimir invariants from the operator algebra]
\label{thm:casimir}
The SU(3)$_c$ Casimir operators take the standard values:
\begin{align}
C_F &= \frac{N^2-1}{2N} = \frac{4}{3}
  \quad (\text{fundamental representation}),
\label{eq:CF}\\
C_A &= N = 3
  \quad (\text{adjoint representation}),
\label{eq:CA}\\
T_F &= \frac{1}{2}
  \quad (\text{fundamental Dynkin index}).
\label{eq:TF}
\end{align}
\end{theorem}

\begin{proof}
From SS\#2 Theorem~1, $T^a = \lambda^a/2$.  Then:
\[
\sum_{a=1}^8 T^a T^a
= \frac{1}{4}\sum_{a=1}^8 \lambda^a\lambda^a
= \frac{4}{3}\,\mathbf{1}_3
\]
by the standard Gell-Mann matrix identity
$\sum_a \lambda^a_{ij}\lambda^a_{kl} = 2(\delta_{il}\delta_{jk}
- \frac{1}{3}\delta_{ij}\delta_{kl})$.
Taking the trace: $\mathrm{Tr}(T^a T^b) = \frac{1}{2}\delta^{ab}$,
giving $T_F = 1/2$.  $C_F = 4/3$ follows from
$\sum_a T^a T^a = C_F \mathbf{1}$.  $C_A = 3$ from
$f^{acd}f^{bcd} = C_A \delta^{ab} = 3\delta^{ab}$.
All values verified numerically to $<10^{-10}$
\citep{cpp_ss_mc}.\qed
\end{proof}

\begin{remark}[Physical significance of the Casimirs]
$C_F = 4/3$ enters the quark self-energy and quark--gluon
coupling in perturbative QCD.  $C_A = 3$ governs
gluon self-interaction (the ``non-Abelian'' contribution).
$T_F = 1/2$ appears in the one-loop $\beta$-function coefficient
$\beta_0 = 11C_A/3 - 4T_F n_f/3 = 11 - 2n_f/3$, which sets
the scale of asymptotic freedom.  The derivation of $\beta_0$
from CPP cage dynamics is the subject of SS\#4.
\end{remark}

%=====================================================================
\section{Non-Abelian Self-Coupling}
%=====================================================================

\subsection{Three-gluon vertex}

In QCD, the Yang-Mills Lagrangian contains a cubic gluon term:
\begin{equation}
\mathcal{L}_{3g}
= -g_s f^{abc}(\partial_\mu A^a_\nu)(A^{b\mu})(A^{c\nu}),
\label{eq:3gluon}
\end{equation}
where $g_s = \sqrt{4\pi\as}$ is the strong coupling.

\begin{theorem}[3-gluon vertex from nested tetrahedral hoppings]
\label{thm:3gluon}
The cubic gluon coupling~\eqref{eq:3gluon} arises in CPP from
the sequential action of two hDP hoppings on the same color state:
\begin{equation}
T^a T^b
= \tfrac{1}{2}\bigl(if^{abc} + d^{abc}\bigr)T^c
+ \tfrac{1}{3}\delta^{ab}\,\mathbf{1},
\label{eq:TaTb}
\end{equation}
where $d^{abc}$ are the symmetric structure constants and
$f^{abc}$ the antisymmetric ones.  The antisymmetric part
$\frac{i}{2}f^{abc}T^c = \frac{1}{2}[T^a,T^b]$ is the
commutator, which gives the non-Abelian cubic coupling.
\end{theorem}

\begin{proof}
Equation~\eqref{eq:TaTb} is the standard decomposition of a
product of SU(3) generators into symmetric and antisymmetric
parts.  The antisymmetric commutator $[T^a,T^b] = if^{abc}T^c$
(SS\#2 Theorem~2 \citep{cpp_ss2}) directly gives the
structure-constant prefactor of~\eqref{eq:3gluon}.\qed
\end{proof}

\begin{remark}[Physical picture]
When a quark at vertex $V_1$ emits gluon $a$ (hopping toward
$V_2$) and the outgoing gluon then interacts with another gluon
$b$ already in the color field, the combined hopping
$T^aT^b - T^bT^a = if^{abc}T^c$ generates a new gluon $c$.
This is the 3-gluon vertex: gluons can scatter off each other
because color is a vertex label and two color changes do not
generally commute.  In QED (photons), $[T^a, T^b] = 0$
(U(1) is Abelian) --- photons do not scatter off each other
at tree level.
\end{remark}

\subsection{Four-gluon vertex}

The quartic gluon coupling $f^{abe}f^{cde}A^aA^bA^cA^d$ follows
from the same mechanism applied twice:
\begin{equation}
[T^a, [T^b, T^c]]
= f^{bce}[T^a, T^e]
= if^{bce}f^{aed}T^d
= -f^{bce}f^{ade}T^d,
\label{eq:4gluon}
\end{equation}
which reproduces the quartic structure constant product
$f^{abe}f^{cde}$ (using the Jacobi identity from SS\#2
Theorem~3 \citep{cpp_ss2}).

%=====================================================================
\section{Color Confinement: Forced Color Neutrality}
%=====================================================================

\begin{theorem}[Color-neutral hadrons from cage geometry]
\label{thm:confinement}
The only stable multi-quark configurations are:
\begin{enumerate}
\item \textbf{Baryons:} three quarks, one at each base vertex
  $\{V_1, V_2, V_3\}$, forming the totally antisymmetric
  color singlet:
  \begin{equation}
  \ket{B} = \frac{1}{\sqrt{6}}\epsilon^{ijk}\ket{i}\ket{j}\ket{k}
  = \frac{1}{\sqrt{6}}(\ket{rgb} - \ket{rbg} + \ket{gbr} - \ket{grb}
    + \ket{brg} - \ket{bgr}).
  \label{eq:baryon_singlet}
  \end{equation}
\item \textbf{Mesons:} quark at vertex $V_i$ and antiquark at
  complementary antivertex $\bar{V}_i$, forming:
  \begin{equation}
  \ket{M} = \frac{1}{\sqrt{3}}\sum_i \ket{i}\ket{\bar{i}}
  = \frac{1}{\sqrt{3}}(\ket{r\bar{r}} + \ket{g\bar{g}}
    + \ket{b\bar{b}}).
  \label{eq:meson_singlet}
  \end{equation}
\end{enumerate}
Both states are annihilated by all eight gluon generators:
$T^a\ket{B} = T^a\ket{M} = 0$ for all $a$.
\end{theorem}

\begin{proof}
For the baryon: $T^a$ acts as a color rotation.  The totally
antisymmetric state $\epsilon^{ijk}\ket{i}\ket{j}\ket{k}$ is
invariant under SU(3) (it is the unique rank-3 antisymmetric
tensor invariant), so $T^a\ket{B} = 0$.

For the meson: $\sum_i\ket{i}\ket{\bar{i}}$ is the SU(3)
invariant contraction of fundamental and antifundamental
representations ($\delta^i_j$), annihilated by all generators.

The cage geometry forces these configurations: a baryon requires
exactly one quark at each of the three base vertices (one quark
per vertex label = one color per vertex = color-neutral).  A
meson requires a quark--antiquark pair on vertex/antivertex
(color/anticolor = color-neutral).  Any other configuration
carries a net color charge and cannot be stabilised by the
tetrahedral cage (the open V$_4$ apex provides no balancing
color).\qed
\end{proof}

%=====================================================================
\section{The QCD Lagrangian from CPP}
%=====================================================================

Combining the results of this series, the full QCD Lagrangian
emerges:

\begin{equation}
\mathcal{L}_{\rm QCD}
= -\frac{1}{4}F^a_{\mu\nu}F^{a\mu\nu}
+ \sum_f \bar{\psi}_f(i\gamma^\mu D_\mu - m_f)\psi_f,
\label{eq:LQCD}
\end{equation}
where:
\begin{itemize}
\item $F^a_{\mu\nu} = \partial_\mu A^a_\nu - \partial_\nu A^a_\mu
  + g_s f^{abc}A^b_\mu A^c_\nu$: field strength from gluon
  edge hoppings (SS\#3, this paper) and structure constants
  (SS\#2 \citep{cpp_ss2}).
\item $D_\mu = \partial_\mu - ig_s T^a A^a_\mu$: covariant
  derivative with Nexus gauge invariance (EW \#5
  \citep{cpp_ew5}, adapted to SU(3)).
\item $m_f$: quark masses from cage binding energy (SS\#1
  \citep{cpp_ss1}; Open Problem~SS-1).
\item Sum over $n_f = 6$ active flavours from SS\#1
  Table~1 \citep{cpp_ss1}.
\end{itemize}

The Yang-Mills coarse-graining limit (EW \#5 Theorem~3
\citep{cpp_ew5}) applies to the strong sector with SU(3)
replacing SU(2): the discrete tetrahedral hopping dynamics
converges to $\mathcal{L}_{\rm QCD}$ as $\Delta s \to 0$.

%=====================================================================
\section{Open Problems}
%=====================================================================

\begin{openprob}[Quark mass formula]
\label{op:mf}
The quark masses $m_f$ in~\eqref{eq:LQCD} are not yet derived
from cage geometry.  The cage-depth argument (SS\#1
Table~1 \citep{cpp_ss1}) gives the qualitative hierarchy
but not the numerical values.
\end{openprob}

\begin{openprob}[Strong coupling $g_s$ from cage geometry]
\label{op:gs}
The strong coupling $g_s = \sqrt{4\pi\as}$ enters all gluon
vertices.  Unlike the EW Weinberg angle (derived to 0.004\%),
$\as(M_Z) \approx 0.118$ has not been derived from 600-cell
geometry.  Strong Series \#4 will address this via the
$\beta$-function.
\end{openprob}

\begin{openprob}[Glueball spectrum]
\label{op:glueball}
Pure gluonic states (glueballs) correspond to closed gluon loops
--- hDP cycles on the tetrahedral cage that do not involve quarks.
The lightest predicted glueball mass from lattice QCD is
$\sim1.5$--$2$~GeV \citep{morningstar1999}.  Deriving this from
CPP tetrahedral loop confinement geometry is an open problem.
\end{openprob}

%=====================================================================
\section{Consolidated Results of SS\#1--3}
%=====================================================================

\begin{table}[H]
\centering
\caption{CPP strong sector: what is derived vs.\ open, after SS\#1--3.}
\label{tab:status}
\begin{tabular}{@{}lll@{}}
\toprule
Result & Status & Paper \\
\midrule
Color = vertex identity & Derived (C$_3$ theorem) & C15, SS\#1 \\
SU(3) algebra $[T^a,T^b]=if^{abc}T^c$ & \textbf{Exact} & SS\#2 \\
$T^a = \lambda^a/2$ & \textbf{Exact} & SS\#2 \\
8 gluon operators identified & \textbf{Exact} & SS\#3 (this) \\
Gluon masslessness & Derived & SS\#3 \\
Gluon spin-1 & Derived & SS\#3 \\
$C_F = 4/3$, $C_A = 3$, $T_F = 1/2$ & \textbf{Exact} & SS\#3 \\
3- and 4-gluon vertices & Derived & SS\#3 \\
Color confinement (geometric) & Derived & SS\#3 \\
Cornell potential $V = -\as/r + \sigma r$ & Reproduced & C14 \\
Asymptotic freedom $\as(Q)$ form & Reproduced & C14 \\
$\beta_0$ from first principles & Open (SS\#4) & --- \\
Quark mass formula & Open (SS\#5) & --- \\
$\as(M_Z) = 0.118$ derived & Open (SS\#4) & --- \\
\bottomrule
\end{tabular}
\end{table}

%=====================================================================
\section{Conclusion}
%=====================================================================

The eight gluons are hDP pairs emitted along the three edges
of the qCP tetrahedral base.  Six color-changing gluons carry
vertex-label transfers ($V_i \to V_j$); two neutral diagonal
gluons carry phase differences.  All eight are massless (open-path
topology, no SS~Vector compression energy), spin-1 (edge orientation
defines polarisation axis), and transform in the adjoint (octet)
representation ($C_A = 3$).

The non-Abelian self-coupling --- the defining feature of
QCD that distinguishes it from QED --- arises from the
antisymmetric part of nested tetrahedral hoppings,
$[T^a, T^b] = if^{abc}T^c$, which was proved exact in
SS\#2 \citep{cpp_ss2}.

Color confinement is geometric: only configurations where all
three base vertices are symmetrically occupied (baryons) or where
vertex/antivertex pairs cancel (mesons) are color-neutral, and
only color-neutral states are stable.

Together, SS\#1--3 complete the derivation of the kinematic
structure of QCD from CPP geometry.  The dynamics
($\beta$-function, quark masses, pion mass) are addressed in
SS\#4--5.

\noindent\textbf{Supplementary material:} Numerical verification
of all Casimir values and commutator algebra at
\texttt{CPP/series\_strong/mc\_su3\_algebra.py}
\citep{cpp_ss_mc}.

%=====================================================================
\bibliographystyle{unsrtnat}
\begin{thebibliography}{14}

\bibitem{cpp_ss1}
T.~L. Abshier and Grok,
``Strong sector overview,''
CPP SS \#1 v1, Hyperphysics Institute (2026).

\bibitem{cpp_ss2}
T.~L. Abshier and Grok,
``SU(3)$_c$ from tetrahedral phase interference,''
CPP SS \#2 v1, Hyperphysics Institute (2026).

\bibitem{cpp_ew1}
T.~L. Abshier and Grok,
``Electroweak sector: introductory overview,''
CPP EW \#1 v3, Hyperphysics Institute (2026).

\bibitem{cpp_ew2}
T.~L. Abshier and Grok,
``The W$^0$ bracelet and W$^\pm$ boson,''
CPP EW \#2 v3.1, Hyperphysics Institute (2026).

\bibitem{cpp_ew5}
T.~L. Abshier and Grok,
``Emergent electroweak unification,''
CPP EW \#5 v3, Hyperphysics Institute (2026).

\bibitem{c14}
T.~L. Abshier and Grok,
``Quark confinement and the Cornell potential from qDP chaining,''
CPP companion~C14, Hyperphysics Institute (2026).

\bibitem{c15}
T.~L. Abshier and Grok,
``Color charge, fractional charge, and SU(3) from tetrahedral cage geometry,''
CPP companion~C15, Hyperphysics Institute (2026).

\bibitem{cpp5014}
T.~L. Abshier and Grok,
``Charge neutrality in baryons and emergent quark charges,''
CPP-5014, Hyperphysics Institute (2026).

\bibitem{cpp_ss_mc}
T.~L. Abshier and Grok,
``CPP strong sector Monte Carlo verification,''
\texttt{mc\_su3\_algebra.py},
CPP/series\_strong, GitHub (2026).

\bibitem{gell-mann1962}
M.~Gell-Mann,
``Symmetries of baryons and mesons,''
\textit{Phys.\ Rev.} \textbf{125}, 1067 (1962).

\bibitem{pdg2026}
Particle Data Group,
``Review of Particle Physics 2026,''
\textit{Phys.\ Rev.\ D} \textbf{110}, 030001 (2026).

\bibitem{gross1973}
D.~J. Gross and F.~Wilczek,
``Ultraviolet behavior of non-Abelian gauge theories,''
\textit{Phys.\ Rev.\ Lett.} \textbf{30}, 1343 (1973).

\bibitem{morningstar1999}
C.~J. Morningstar and M.~J. Peardon,
``Analytic smearing of SU(3) link variables in lattice QCD,''
\textit{Phys.\ Rev.\ D} \textbf{60}, 034509 (1999).

\bibitem{abshier_sr}
T.~L. Abshier,
``Mechanistic derivation of relativistic effects via Space Stress Vector,''
viXra:2511.0062 (2025).

\end{thebibliography}

\end{document}
